\documentclass[10pt,a4paper]{amsart}
\usepackage[margin=19mm]{geometry}
\usepackage{amsmath,amssymb,mathtools}
\usepackage[scr=rsfs,cal=euler]{mathalfa}
\usepackage{xfrac}
\usepackage[unicode,hidelinks,bookmarks=false]{hyperref}
\newcommand{\ie}{i.e.\ }
\newcommand{\cf}{cf.\ }
\newcommand{\resp}{respectively\ }
\newcommand{\Subsection}{\subsection}
\providecommand{\nobreakdash}{\nobreak\hbox{-}}
\setlength{\emergencystretch}{2em}
\multlinegap0pt
\usepackage{bm}
\usepackage{bbm}
\usepackage{dsfont}
\usepackage{xspace}
\usepackage{cancel}
\usepackage{mathtools}
\usepackage{amsthm}
\makeatletter
\@ifundefined{mydef}{\newtheorem{mydef}{Mydef}}{}
\@ifundefined{myclaim}{\newtheorem{myclaim}{Claim}}{}
\@ifundefined{mylem}{\newtheorem{mylem}[mydef]{Lemma}}{}
\makeatother
\DeclareMathOperator{\ISNIC}{ISNIC}
\newcommand{\dQ}{\mathbb{Q}}
\title[Disjoint Stationary Sequences on an Interval of Cardinals]{Disjoint Stationary Sequences on an Interval of Cardinals}
\author{Hannes Jakob}
\date{Source: arXiv:2309.01986v2 (CC BY 4.0)}
\begin{document}
\maketitle

\noindent\emph{Source and licence.} Disjoint Stationary Sequences on an Interval of Cardinals. Version arXiv:2309.01986v2 (CC BY 4.0), distributed under the Creative Commons Attribution 4.0 International licence, as recorded in the arXiv licence metadata for this exact version and shown on the abstract page https://arxiv.org/abs/2309.01986v2. Full rights notice: licenses/arxiv-2309.01986-CC-BY-4.0.txt.\par\smallskip

\noindent\textit{Editorial scope.} Counted unit: the Mylem whose source label is \texttt{PreserveDown} in the article (source0.tex lines 757--759, source label \texttt{PreserveDown}), delivered together with the article's own complete original proof of it (source0.tex lines 761--793, one \texttt{proof} environment of 33 lines). No mathematics is written, repaired or re-derived: statement and proof are the article's own text line for line. Required context retained uncounted: none. Every definition, display and label of those lines is retained unchanged. Omitted from this package and not redistributed: 1--756, 760--760, 794--948. Nothing inside a retained excerpt is omitted or altered: every retained body is delivered exactly as the article's lines are written, so no omission receipt of any kind accompanies this package. Citations: the retained excerpts carry no citation command at all, so no bibliography entry of the article is transcribed and no reference is redistributed. Reference adaptation: none was needed and none was made; every reference inside the delivered unit resolves inside it, so no unresolved reference or citation is produced. Printed slips kept literal: no printed word, symbol, hypothesis or number of the counted statement or of its proof was altered. Layout and class: the article's own class is \texttt{amsart} with packages including babel, amssymb, mathrsfs, hyperref, xcolor, eucal, amsmath, amsthm, imakeidx, breqn, diagbox, verbatim, fancyhdr, tikz-cd; the wrapper is the lane's pinned \texttt{amsart} editorial wrapper at 10pt on A4, which restates the theorem-like environments the retained text needs and adds the article's own macro definitions that the retained text needs (\texttt{\textbackslash ISNIC}, \texttt{\textbackslash dQ}), with extra wrapper packages (bm, bbm, dsfont, xspace, cancel, mathtools). The delivered numbering is the unit's own. Assets: no retained span contains a figure environment, a graphics-inclusion command, a PDF-page inclusion, a table or a TikZ picture; no image file is copied into this package, the package declares no asset dependency and no image file is redistributed with it. Licence: version arXiv:2309.01986v2 of this article is distributed under the Creative Commons Attribution 4.0 International licence, as recorded in the arXiv licence metadata for this exact version and shown on the abstract page https://arxiv.org/abs/2309.01986v2. A full rights notice is delivered at licenses/arxiv-2309.01986-CC-BY-4.0.txt. Characters: every retained excerpt is pure ASCII, so the delivered engine, XeLaTeX with its default Latin Modern text font, reports no missing character.
	\begin{mylem}\label{PreserveDown}
		Let $\mu>\aleph_0$ be a regular cardinal and $\dQ$ a poset. Assume that $\ISNIC^+(\mu^+)$ holds in $V[G]$, where $G$ is $\dQ$-generic, and $\dQ$ is ${<}\,\mu^+$-distributive. Then $\ISNIC^+(\mu^+)$ holds in $V$.
	\end{mylem}

	\begin{proof}
		In $V$, let $\Theta$ be arbitrary and let $F\colon[H^V(\Theta)]^{<\omega}\to[H^V(\Theta)]^{\mu}$ be any function. Let $\Theta'$ be so large that $F,\dQ\in H(\Theta')$. The following is clear:
		\setcounter{myclaim}{0}
		\begin{myclaim}
			In $V[G]$, the set
			$$D:=\{M\in[H^V(\Theta')]^{\mu}\;|\;M[G]\cap V=M\}$$
			is club in $[H^V(\Theta')]^{\mu}$.
		\end{myclaim}
		
		Since $H^V(\Theta')[G]=H^{V[G]}(\Theta')$ by the assumption $\dQ\in H(\Theta')$, we know that $\{M[G]\;|\;M\in D\}$ is club in $[H^{V[G]}(\Theta')]^{\mu}$. So by $\ISNIC^+(\mu^+)$ holding in $V[G]$, we can find $M\in D$ such that $F\in M$, $H^V(\Theta)\in M$, $M[G]\prec(H^{V[G]}(\Theta'),\in)$ and $M[G]$ is internally stationary but does not contain a club in $[M[G]\cap\mu^+]^{<\mu}$. Clearly $N:=M[G]\cap H^V(\Theta)=M\cap H^V(\Theta)$ is closed under $F$, since $F\in V$. Also $N\in V$ by the distributivity of $\dQ$.
		
		\begin{myclaim}
			$N$ is internally stationary in $V$.
		\end{myclaim}
		
		\begin{proof}
			In $V$, let $c\subseteq[N]^{<\mu}$ be club. Then $\{m\in[M]^{<\mu}\;|\;m\cap H^V(\Theta)\in c\}$ is club in $[M]^{<\mu}$ and so $\{m[G]\;|\;m\in[M]^{<\mu}\wedge m\cap H^V(\Theta)\in c\}$ is club in $[M[G]]^{<\mu}$. Ergo there is $m\in[M]^{<\mu}$ with $m\cap H^V(\Theta)\in c$ such that $m[G]\in M[G]$ since $M[G]$ is internally stationary. As in the first claim, we may assume $m[G]\cap V=m$ from which it follows that
			$$m\cap H^V(\Theta)=m[G]\cap H^V(\Theta)\in c\cap M[G]\cap H^V(\Theta)=c\cap N$$
			since $m\cap H^V(\Theta)\in H^V(\Theta)$ by its size and the distributivity of $\dQ$.
		\end{proof}
		
		We are finished after showing:
		
		\begin{myclaim}
			$N$ does not contain a club in $[N\cap\mu^+]^{<\mu}$ in $V$.
		\end{myclaim}
		
		\begin{proof}
			Let $c\subseteq[N\cap\mu^+]^{<\mu}$ be club. Since $N\cap\mu^+=M[G]\cap\mu^+$, $c\subseteq N$ would imply that $M[G]\supseteq N$ contains a club in $[M[G]\cap\mu^+]^{<\mu}$, a clear contradiction.
		\end{proof}
		
		This finishes the proof of Lemma \ref{PreserveDown}.
	\end{proof}

\end{document}
